% ============================================================
% ประวัติผู้เรียบเรียง (Author Biography)
% ============================================================
\chapter*{ประวัติผู้เรียบเรียง}
\addcontentsline{toc}{chapter}{ประวัติผู้เรียบเรียง}

\noindent\textbf{ผู้ช่วยศาสตราจารย์ ดร.ชีวะ ทัศนา (Asst. Prof. Dr. Chewa Thassana)}

\vspace{0.4cm}
\noindent\textbf{ตำแหน่งทางวิชาการและสังกัด}
\begin{itemize}
    \item ผู้ช่วยศาสตราจารย์ ประจำสาขาวิชาฟิสิกส์และวิทยาศาสตร์ทั่วไป
    \item หัวหน้ากลุ่มวิจัยฟิสิกส์เกษตรดิจิทัลและปัญญาประดิษฐ์ฝังตัว (Digital Agriphysics and Embedded AI Research Group)
    \item คณะวิทยาศาสตร์และเทคโนโลยี มหาวิทยาลัยราชภัฏรำไพพรรณี
\end{itemize}

\vspace{0.3cm}
\noindent\textbf{ประวัติการศึกษา}
\begin{itemize}
    \item ปริญญาเอก \quad วิศวกรรมศาสตรดุษฎีบัณฑิต (วศ.ด.) สาขาวิชาวิศวกรรมไฟฟ้าและอิเล็กทรอนิกส์
    \item ปริญญาโท \quad วิศวกรรมศาสตรมหาบัณฑิต (วศ.ม.) สาขาวิชาวิศวกรรมโทรคมนาคม
    \item ปริญญาตรี \quad วิทยาศาสตรบัณฑิต (วท.บ.) สาขาวิชาฟิสิกส์ประยุกต์
\end{itemize}

\vspace{0.3cm}
\noindent\textbf{ความเชี่ยวชาญทางวิชาการและงานวิจัย}
\begin{itemize}
    \item ระบบสมองกลฝังตัวและอินเทอร์เน็ตของสรรพสิ่งสำหรับการเกษตรแม่นยำสูง (AIoT in Precision Agriculture)
    \item โครงข่ายเซนเซอร์ไร้สายพลังงานต่ำและระบบสื่อสารแบบเมช (Zigbee 3.0, LoRaWAN, BLE Mesh)
    \item การประมวลผลสัญญาณฟิสิกส์และปัญญาประดิษฐ์ที่ขอบโครงข่าย (Edge Computing and TinyML)
    \item การออกแบบและพัฒนาระบบอัตโนมัติในสวนผลไม้มูลค่าสูงและการเพาะเลี้ยงสัตว์น้ำชายฝั่ง
\end{itemize}

\clearpage
